% !TeX program = xelatex
\documentclass[12pt]{nextbook}

\UseTemplateSet{
  globalfonts = {libertinus, shanggu},
  features = {hyperlinks}
}

\newcommand{\ReaderColumnSep}{4mm}
\newcommand{\ReaderVocabSeparator}{}
\usepackage{paracol}
\usepackage{xcolor}
\usepackage{eso-pic}
\usepackage{needspace}

\vbadness=10000
\hbadness=10000
\hfuzz=3pt
\emergencystretch=1.5em

\makeatletter
\AtBeginDocument{%
  \@ifpackageloaded{hyperref}{%
    \pdfstringdefDisableCommands{%
      \def\ruby#1#2{#1}%
      \def\ReaderVocabText#1{#1}%
      \def\ReaderVocabTerm#1{#1}%
    }%
  }{}%
}
\makeatother

\providecommand{\ReaderColumnSep}{6mm}
\providecommand{\ReaderColumnRatio}{0.75}
\providecommand{\ReaderVocabBaselineskip}{20pt}
\providecommand{\ReaderVocabText}[1]{\textbf{#1}}
\providecommand{\ReaderVocabTerm}[1]{\textbf{#1}}
\providecommand{\ReaderVocabSeparator}{:}
\providecommand{\ReaderVocabItem}[2]{\noindent\ReaderVocabTerm{#1}\ReaderVocabSeparator\ #2\par}

\setlength{\columnsep}{\ReaderColumnSep}
\setlength{\columnseprule}{0pt}
\columnratio{\ReaderColumnRatio}
\swapcolumninevenpages

\newcounter{readersection}
\newwrite\vocabfile
\newif\iffirstreaderlabelwritten
\newlength{\readercolumnswidth}
\newlength{\readertitlegutter}
\setlength{\readercolumnswidth}{\dimexpr\textwidth-\columnsep\relax}
\setlength{\readertitlegutter}{.75\columnsep}

\makeatletter
\newcommand{\readerseparatorline}{%
  \if@mainmatter
    \AtTextLowerLeft{%
      \ifodd\value{page}%
        \put(\LenToUnit{\dimexpr .75\readercolumnswidth+.5\columnsep\relax},0){%
          \color{black!35}\rule{0.3pt}{\textheight}%
        }%
      \else
        \put(\LenToUnit{\dimexpr .25\readercolumnswidth+.5\columnsep\relax},0){%
          \color{black!35}\rule{0.3pt}{\textheight}%
        }%
      \fi
    }%
  \fi
}
\makeatother

\AddToShipoutPictureBG{\readerseparatorline}

\newcommand{\currentvocabfile}{\jobname-vocab-\thereadersection.voc}
\newcommand{\beginvocabwrite}{\immediate\openout\vocabfile=\currentvocabfile}
\newcommand{\finishvocabwrite}{\immediate\closeout\vocabfile}
\newcommand{\vocabitem}[2]{\ReaderVocabItem{#1}{#2}}
\newcommand{\printvocablist}{%
  \footnotesize
  \setlength{\parindent}{0pt}%
  \setlength{\parskip}{0pt}%
  \setlength{\baselineskip}{\ReaderVocabBaselineskip}%
  \raggedright
  \input{\currentvocabfile}%
}

\newcommand{\readersectiontitle}[1]{%
  \par
  \begingroup
  \ifodd\value{page}%
    \leftskip=\readertitlegutter%
    \rightskip=0.25\textwidth plus 1fil%
  \else
    \leftskip=\dimexpr .25\readercolumnswidth+.5\columnsep+\readertitlegutter\relax%
    \rightskip=0pt plus 1fil%
  \fi
  \section{#1}%
  \par
  \endgroup
}

\NewDocumentCommand{\voc}{omm}{%
  \ReaderVocabText{#2}%
  \IfNoValueTF{#1}{%
    \immediate\write\vocabfile{\string\vocabitem{#2}{#3}}%
  }{%
    \immediate\write\vocabfile{\string\vocabitem{#1}{#3}}%
  }%
}

\newenvironment{readersection}[1]{%
  \stepcounter{readersection}%
  \beginvocabwrite
  \Needspace{10\baselineskip}%
  \iffirstreaderlabelwritten\else
    \label{reader:firstpage}%
    \global\firstreaderlabelwrittentrue
  \fi
  \readersectiontitle{#1}%
  \begin{paracol}{2}
}{%
  \finishvocabwrite
  \switchcolumn
  \printvocablist
  \end{paracol}
}


\title{English Vocabulary Builder}
\author{Kelvin Yang}
\date{\today}

\begin{document}
\frontmatter
\maketitle
\tableofcontents
\mainmatter
\chapter{4k-Level}
\begin{readersection}{The Museum Project}
When Ms. Hart asked us to design a new museum exhibition, she began with a \voc{questionnaire}{問卷}. She wanted to know what visitors liked: science, history, music, politics, animals, or old machines. The answers were placed on a metal \voc{rack}{架子} beside the door, where anyone could read them. At first, I did not want to \voc{dwell}{居住} on the details, but soon I saw that every answer revealed a small world.

The exhibition was called “A Small \voc{realm}{領域} of Human Life.” One room showed the crown of a forgotten \voc{monarch}{君主}, together with notes from the old \voc{legislature}{立法機關}. Another room was about \voc{astronomy}{天文學}; a child with a bandaged \voc{ankle}{腳踝} stood under the painted stars and asked why the moon seemed to \voc{rotate}{旋轉} around us. Our school had a fixed \voc{quota}{限額} of visitors, so everyone had to \voc{conform}{遵從} to the timetable and \voc{comply}{遵守} with the safety rules.

The museum itself had a strange history. Its \voc{predecessor}{前身} was a public library, and before that it had been a poor \voc[dwell]{dwelling}{居住} for workers who cut \voc{timber}{木材}. A photograph near the entrance was so old that the faces had begun to \voc{blur}{模糊}; still, you could see a man’s tired \voc{forehead}{前額}, a woman holding a paper \voc{packet}{小包}, and a child standing near a \voc{bundle}{一捆} of books tied with a \voc{cord}{繩子}. On a small label, someone had written in blue \voc{ink}{墨水}: “Do not \voc{discard}{丟棄} memory too quickly.”

Our guide, Mr. Lane, was \voc{fond}{喜愛的} of strange connections. In the agricultural room, he spoke about \voc{olive}{橄欖} trees, \voc{mammal}{哺乳動物} bones found near an old \voc{barn}{穀倉}, and the \voc{beneficial}{有益的} effects of teaching children to \voc{cultivate}{培養} a garden. He also explained how people learned to \voc{recycle}{回收} metal, glass, and \voc{aluminium}{鋁}, and how old \voc{wax}{蠟} tablets could preserve writing. He said that even a broken \voc{fork}{叉子}, a torn \voc{sleeve}{袖子}, or a piece of \voc{textile}{紡織品} could help historians understand daily life.

In the political room, the story became darker. There had once been a violent \voc{faction}{派系} that spread a \voc{conspiracy}{陰謀} about the royal family. Some people were \voc[fine]{fined}{罰款}, and others were forced into \voc{exile}{流亡}. A serious \voc{sergeant}{軍士} stood in one photograph, guarding a box of \voc{ballot}{選票} papers. Nearby, a judge’s notes described a harsh \voc[sentence]{sentencing}{判刑} hearing, a case of \voc{torture}{酷刑}, and the attempt of one prisoner to \voc{regain}{重獲} his dignity. The museum did not try to \voc{obscure}{掩蓋} these events. Instead, it showed how fear can become \voc{toxic}{有害的}, how a single \voc{bullet}{子彈} can change a city, and how no government can \voc{foresee}{預見} every \voc{obstacle}{障礙} that history will place before it.

Still, the exhibition was not hopeless. The next room was about medicine. A famous \voc{surgeon}{外科醫生} had once worked there after a war. On display were his \voc{hardware}{器械} tools, a \voc{horizontal}{水平的} table, a \voc{vertical}{垂直的} chart, and a diagram showing how \voc{oxygen}{氧氣} moves through the body. A model of a lung showed tiny \voc[particle]{particles}{微粒} in the air. The surgeon believed that illness could \voc{impair}{損害} the body, but kindness could heal the mind. He used to say that survival was not a \voc{miracle}{奇蹟}; it was the result of skill, care, and patience.

The science room had a more playful tone. It compared \voc{physics}{物理學} with daily experience. One machine used \voc{copper}{銅} wire and a glass \voc{dome}{圓頂}; another machine tried to \voc{align}{對齊} light through a small hole. A child asked whether a \voc{hybrid}{混合型的} car belonged in the science room or the social room. Mr. Lane laughed and said that this was the museum’s central \voc{paradox}{悖論}: every object belonged to more than one category. Even a simple map could become a \voc{metaphor}{隱喻} for human choices.

At noon, we ate in the museum café. It was designed to \voc{cater}{迎合} to students, elderly people, and tourists. The food was not \voc{premium}{高級的}, but it was warm. A singer from the local \voc{choir}{合唱團} performed a short \voc{solo}{獨唱}; later, a small \voc{jazz}{爵士樂} group played near the window. A little girl sat on her father’s \voc{lap}{膝上}, watching the sunlight \voc{gleam}{閃光} on a marble statue. The \voc{bride}{新娘} in a wedding photograph wore a thin \voc{veil}{面紗}; on the \voc{eve}{前夕} of her marriage, she had donated her family letters to the museum.

In the technology room, there were old \voc{telecom}{電信} devices, a wooden phone box, and early computer parts. A sign explained that digital \voc{literacy}{素養} was now as important as reading paper books. Some devices looked \voc{dim}{昏暗的} and ugly, but visitors were \voc{ambitious}{有抱負的} enough to try them. One screen suddenly made a loud \voc{jerk}{急動}, then went black. A boy almost took a \voc{tumble}{跌倒}, but he caught the edge of a table. The staff member in charge had to \voc{administer}{施行} a quick check and then put a warning sign near the \voc{exit}{出口}.

In a side gallery, the museum showed how art and society often overlap. One poster discussed the history of \voc{homosexual}{同性戀的} identity and the struggle for public recognition. Another exhibit asked why some societies \voc{classify}{分類} people too quickly, turning living individuals into fixed types. Mr. Lane said that a museum should not \voc{preoccupy}{佔據心思} visitors with labels alone. “A label can help,” he said, “but it can also become a \voc{foul}{惡劣的} habit if it replaces thought.” We took a slow \voc{stroll}{散步} through the gallery, and I began to understand that the \voc{pursuit}{追求} of knowledge was not the same as the possession of it.

The last room was the most \voc{spectacular}{壯觀的}. It looked like a stone \voc{temple}{神殿}, with a white \voc{marble}{大理石} floor and a painted sky above. Visitors could \voc{bathe}{沐浴} their hands in a shallow basin of light, not water. A sign asked them to place each object into one group: tool, artwork, weapon, \voc{commodity}{商品}, or sacred object. Some people argued. One man called the task an \voc{insult}{侮辱}; another said that a museum should never treat a sacred image like a \voc{beast}{野獸} in a cage. Mr. Lane quietly said that disagreement was part of the public \voc{arena}{公共舞臺}.

Before we left, the museum held a small charity \voc{auction}{拍賣}. Each \voc{recipient}{接受者} of a donated object had written a short note about why it mattered. One case contained a broken toy, a cracked cup, and a note about a manufacturing \voc{defect}{缺陷}. The objects were not famous; they had no \voc{fame}{名聲} at all. Yet the label said that every ordinary object had an \voc{eligible}{有資格的} place in history if it helped someone understand the past. The museum was \voc[affiliate]{affiliated}{使附屬} with a university, but it did not speak only to experts. It asked ordinary visitors to look carefully, think slowly, and remember that the past is never only one story.
\end{readersection}

\begin{readersection}{The Cave Archive}

When the old railway museum received a bundle of papers from a mountain town, the director hoped the files would arrive \voc{intact}{完整的}. They did not. A yellow \voc{fax}{傳真} from twenty years before lay on top, but the ink had blurred, and a rusted \voc{spike}{尖釘} from the abandoned track had torn a corner from a handwritten \voc{verse}{詩句}. The author was \voc{anonymous}{匿名的}; he used old pronouns like \voc{thou}{你} and wrote some nouns in strange \voc{plural}{複數} forms. Beside the papers was a dry \voc{mushroom}{蘑菇}, a \voc{crude}{粗糙的} map, and a small \voc{cartoon}{漫畫} of a man with a huge \voc{chin}{下巴}. It looked \voc{comic}{滑稽的}, but Dr. Lane, an \voc{amateur}{業餘的} historian with serious habits, refused to laugh.

The story led us to a \voc{cave}{洞穴} above the town. Local memory said that the \voc{predominant}{主要的} use of the cave had once been storage, but no clear \voc{classification}{分類} existed for the objects inside. A hunter had seen \voc{deer}{鹿} near the entrance, and the first \voc{scenario}{情景} we imagined was simple: smugglers had hidden commodities there. Yet each new clue seemed to \voc{diminish}{減少} that theory. The museum had to \voc{repay}{償還} a small grant if the project failed, so we tried to \voc{fuse}{融合} careful research with practical speed. Our only \voc{trophy}{獎品} would be knowledge.

The town council agreed to \voc{furnish}{提供} lamps, \voc[mat]{mats}{墊子}, and basic hardware. The first \voc{sensation}{感覺} inside the cave was cold air mixed with the \voc{odour}{氣味} of damp stone and old \voc{ash}{灰}. A child from the school choir watched a dog \voc{lick}{舔} water from a shallow pool, then \voc{stagger}{踉蹌} away when a bat rushed past his forehead. We planned to \voc{excavate}{挖掘} only a narrow area, because eggs might still \voc{hatch}{孵化} in nests near the roof. Nobody wanted his \voc{ego}{自尊} to become bigger than the cave itself.

The university \voc{chancellor}{校長} visited on the second morning. She stood beside a broken \voc{statue}{雕像} and asked about the \voc{duration}{持續時間} of the work. I showed her a \voc{diagram}{圖表} in which vertical shafts \voc{complement}{補充} horizontal passages. On a folding table stood cheap \voc{champagne}{香檳}, a rubber mat, a \voc{hammer}{錘子}, and a fork someone had used for lunch. We tried to spread the lights \voc{evenly}{均勻地}, but one bulb began to \voc{explode}{爆裂} with tiny sparks, so a sergeant from the safety team cut the cord and replaced it.

At the back of the cave we found a rack of manuscripts and a small packet of letters. Their writer had a deep \voc{longing}{渴望} to regain fame. He had drawn a cartoon showing a medieval engineer seeking a \voc{patent}{專利} for a rotating machine. The papers also contained a \voc{medieval}{中世紀的} \voc{array}{一系列} of symbols: a crown, a \voc{pope}{教宗}, a \voc{butterfly}{蝴蝶}, a \voc{bolt}{螺栓}, and a \voc{spoon}{湯匙}. Marta raised one \voc{eyebrow}{眉毛}; the pattern seemed obscure. Dust made her \voc{choke}{嗆住}, and she said the cave should not \voc{deprive}{剝奪} anyone of oxygen for the sake of an old metaphor.

My task was \voc{linguistic}{語言的}. I compared the strange verse with temple inscriptions and with a questionnaire from a literacy project. The language showed the \voc{legacy}{遺產} of several realms: a monarch’s court, a \voc{socialist}{社會主義的} printing club, and a church choir. Some translations were not \voc{satisfactory}{令人滿意的}, because the writer loved to \voc{boast}{吹噓} and to conceal ordinary facts under grand words. From a \voc{ridge}{山脊} above the town, he compared himself to an \voc{emperor}{皇帝}; in a \voc{comparative}{比較的} note, he called his predecessor only a beast. That insult told us more about his character than about history.

The third day brought danger. A student slipped on marble dust and got a \voc{bruise}{瘀傷} on his ankle. Near him, a wet wall held a living \voc{organism}{生物} that no one could classify. To preserve scientific \voc{integrity}{誠信}, we refused to discard it as mere mould. We made a \voc{stack}{堆} of samples and labelled them carefully. One assistant began to \voc{obsess}{著迷} over a possible conspiracy, claiming that the writer had tried to \voc{betray}{背叛} his faction. Dr. Lane told him to sleep \voc{overnight}{過夜} before making any inference.

By morning the cave had become an arena of debate. The bus \voc{fare}{車費} for students was low, yet the town’s \voc{infrastructure}{基礎設施} was weak, and the road was almost impassable after rain. Someone found a manuscript about an old pope, written beside a drawing of a butterfly. Another page gave a \voc{grim}{嚴峻的} warning: “If thou break this pledge, the mountain will fall.” Beneath it was a \voc{signature}{簽名}, still clear in copper-coloured ink. The valley had \voc{abundant}{豐富的} timber, but its people had made a \voc{pledge}{承諾} not to cut the forest around the cave.

The council did not want to remain \voc{passive}{被動的}. They asked whether the warning was \voc{applicable}{適用的} to modern law, and whether the museum could turn the site into a \voc{viable}{可行的} education centre. We replaced each weak \voc{bulb}{燈泡}, tightened every bolt, and marked the exit. At supper, an accountant joked that even a spoon from the cave would be put in an auction. The mayor disliked the joke because some officials wanted to \voc{privatise}{私有化} the land. A night \voc{patrol}{巡邏} was arranged to prevent theft.

That conflict \voc[compel]{compelled}{迫使} us to speak openly. If the town tried to \voc{conceal}{隱瞞} the finds, the loss would be of great \voc{magnitude}{重大性}. The chancellor offered to \voc{endow}{資助} an \voc{interim}{臨時的} archive, but warned us not to let a \voc{con}{騙局} grow around local pride. Her \voc{outlook}{觀點} was practical: begin with a \voc{shallow}{淺的} excavation, then \voc{opt}{選擇} for deeper work only after a full report. Her words lit a \voc{spark}{火花} in the room, though not everyone was pleased.

The doctor from the clinic added another issue. Dust in the cave could cause \voc{acute}{急性的} breathing problems, and one volunteer had \voc{diabetes}{糖尿病}, so long hours without food were unsafe. We \voc[refine]{refined}{完善} the schedule into small \voc[chunk]{chunks}{塊} of time and wrote elementary rules on the wall: drink water, leave when dizzy, and never work alone near the \voc{dock}{碼頭} by the underground pool. A \voc{rap}{說唱} group from the school offered to make a song about the project, but the teacher feared that rap would make the discovery seem too sensational. The students argued that music could help every \voc{constituency}{選民群體} in town feel included.

In the final report, I tried to \voc{infer}{推斷} where the whole tradition might \voc{originate}{起源}. The answer still \voc[fascinate]{fascinates}{吸引} me. A medieval craftsman may have hidden his designs there, and later readers added verse, political notes, and religious symbols. We had to refine this theory again and again, because a single defect could impair the whole argument. Yet the conclusion was strong enough for an \voc{elementary}{初級的} exhibition: a cave once used for storage became a place where poetry, technology, faith, and local politics met. Its legacy was not a miracle, but something more useful: a patient record of how ordinary people tried to preserve their world.

\end{readersection}

\begin{readersection}{The Town Inquiry}
On the edge of North \voc{Avenue}{大道} stood a broken research house that everyone simply called the old \voc{greenhouse}{溫室}. For years it had been locked by a suspicious \voc{landlord}{房東}, and rumours about \voc{fraud}{欺詐}, strange experiments, and missing public money moved through the town like smoke. When \voc{Ms}{女士} Hart, our history teacher, asked our class to prepare an \voc{anniversary}{周年紀念} exhibition about the place, she did not want a dramatic \voc{poster}{海報} or a comic \voc{script}{腳本}. She wanted evidence, order, and \voc{rigour}{嚴格}.

Our \voc{preliminary}{初步的} visit began in the yard, where \voc{clay}{黏土} pots lay under a rusted rack and a wooden \voc{mould}{模具} had grown dark with damp. A chemistry table still held a label for \voc{sodium}{鈉} and another for carbon \voc{dioxide}{二氧化物}, while a cracked \voc[simulate]{simulator}{模擬} stood beside a model \voc{rocket}{火箭}. Its \voc{terminal}{終端} screen was dead, but a careless touch made its engine \voc{vibrate}{震動}, and we all stepped back. “Do not \voc{stab}{刺} at the buttons,” Ms Hart said, pointing at a narrow \voc{slot}{狹槽} in the panel. We had to \voc{scramble}{匆忙行動} for notebooks because the machine began to make a \voc{peculiar}{奇特的} \voc{murmur}{低語}, as if it still remembered its old work.

Inside, a \voc{ladder}{梯子} led to a storage room above the laboratory. We climbed slowly, mindful of every \voc{wrist}{手腕} and ankle, because the steps were \voc[batter]{battered}{連續擊打} and the floor seemed ready to \voc{crumble}{碎裂}. \voc{Amid}{在……之中} the dust, we found \voc{armour}{盔甲} from a school play, a \voc{primitive}{原始的} \voc[optic]{optical}{視覺的；光學的} device, and a box of \voc[needle]{needles}{針} used in plant experiments. There was also a \voc{confidential}{機密的} file labelled “\voc{abnormal}{異常的} growth.” It described how one scientist tried to simulate disease in leaves, then claimed that a \voc{clone}{克隆} of a rare plant could save the town’s economy. The idea had \voc{relevance}{相關性}, but the method seemed careless.

Ms Hart asked me to \voc{compile}{彙編} a list of objects while Lena took photographs. I wrote down clay, mould, sodium, dioxide, and \voc{mineral}{礦物} samples; then I noted \voc{grease}{油脂} on the old \voc{keyboard}{鍵盤} and a sharp \voc{lemon}{檸檬} smell from a sealed bottle. A \voc{physician}{醫生}’s note warned that one worker’s \voc{liver}{肝臟} had begun to \voc{deteriorate}{惡化} after \voc[prolong]{prolonged}{延長} exposure to chemicals. The note was not a final diagnosis, but its tone was \voc{stern}{嚴厲的}. It also said that alcohol, especially \voc{whiskey}{威士忌}, had been found in the office, though nobody knew whether it belonged to the workers or to the landlord.

The town archive helped us understand the wider story. A \voc{capitalist}{資本主義的；資本家} company had once tried to build a \voc{monopoly}{壟斷} over a new seed variety. Its managers \voc[aspire]{aspired}{渴望} to fame and spoke of a \voc{mighty}{強大的} future, but their \voc{hierarchy}{等級制度} was brutal. Workers \voc[march]{marched}{行進} through the avenue in protest, saying the company would \voc{displace}{取代；迫使離開} small farmers and \voc{oppress}{壓迫} anyone who refused to comply. Some residents \voc[deem]{deemed}{認爲} the protest unnecessary; others saw it as a \voc{sincere}{真誠的} attempt to defend \voc{ecological}{生態的} life. The debate gained \voc{momentum}{勢頭}, and the company’s public image began to deteriorate.

The most moving document was a letter from an old gardener. He did not write like a scholar, yet the \voc{essence}{本質} of his argument was clear: “A plant is not a \voc{proposition}{命題} in a \voc{thesis}{論題；論文}; it lives in soil, water, light, and care.” He asked readers to \voc{contemplate}{沉思} whether progress without justice could ever be called progress. He also warned against \voc{prejudice}{偏見} toward poor labourers, because some officials had called them \voc{idle}{閒置的；懶散的}, \voc{arrogant}{傲慢的}, or even used the offensive word \voc{bastard}{私生子；混蛋}. Ms Hart \voc[frown]{frowned}{皺眉} when she read that line, and said such language was meant to \voc{humiliate}{羞辱} people.

At \voc{supper}{晚餐} that evening, our group met in a small café near the dock. The owner had \voc[brew]{brewed}{沖泡；釀造} tea with \voc{spice}{香料} and lemon, and the warm smell helped us relax. We discussed the \voc{dual}{雙重的} purpose of our exhibition: it had to tell a \voc{straightforward}{直截了當的} story, but it also had to show the \voc{infinite}{無限的} complexity behind one public scandal. Tom wanted to \voc{omit}{省略} the painful parts, especially the \voc{fatal}{致命的} accident in which a worker fell from a roof during a \voc{blaze}{火焰}. I disagreed. To \voc{deter}{阻止} future negligence, we had to include what people would rather forget.

The accident report was hard to read. It said the fire \voc[erupt]{erupted}{爆發} after a lamp fell into grease, and the blaze spread through dry timber. A worker tried to \voc{despatch}{派遣；發送} a warning by fax, but the line failed. Others were \voc[detain]{detained}{拘留；扣留} by locked doors, and one man survived only because a \voc{veterinarian}{獸醫} from the next street broke a window and pulled him out. The fire destroyed the laboratory but did not \voc{wreck}{毀壞} every record. Some papers were \voc[box]{boxed}{裝箱；拳擊} in metal cases; others were saved by a woman who later became a local councillor.

We also discovered a strange cultural layer. A small \voc{cult}{崇拜；邪教} of admirers had formed around the chief scientist after his death. They treated his notes like a \voc{gospel}{福音}, repeating every claim without question. In one notebook he had written, “The old world must be remade.” This sounded \voc{superb}{極好的} to some readers, but to us it seemed dangerous. It could \voc{obstruct}{阻礙} ordinary judgement. Ms Hart said that a slogan can \voc{illuminate}{照亮；闡明} an age, but it can also conceal the damage done by those who think their own vision is pure.

Near the end of the project, we invited the town’s oldest resident to speak. Her \voc{ancestor}{祖先} had worked in the greenhouse. She remembered stories of \voc[owl]{owls}{貓頭鷹} in the roof, of workers who boxed in the courtyard after work, and of an engineer who could not \voc{blink}{眨眼} when he was angry. She laughed at some memories, then became quiet. “People thought the case was simple,” she said. “It was not. There was hope, fear, money, science, and human pride. That is why you must not make the exhibition too \voc{minimal}{最小的}.”

Her words shaped our final \voc{stance}{立場}. We placed the clay pots near the chemical labels, the worker’s letter beside the company contract, and the accident report under a dim lamp. We used \voc[stitch]{stitches}{縫線；縫合} of red thread to connect photographs, maps, and testimony. The thread looked almost like a \voc{pulse}{脈搏} running through the display. Above it, our main poster asked one question: “What happens when ambition loses its care for people?”

On opening day, the exhibition hall was full. Some visitors were \voc{furious}{憤怒的} because we had not praised the company’s achievements enough. Others said the display was honest. The landlord stood at the back, silent, while the old greenhouse appeared on the screen behind him. For a moment nobody spoke. Then Ms Hart gave a small nod. Our work had not solved every problem, but it had made the town \voc{comprehend}{理解} why the past still \voc[loom]{loomed}{隱約出現} over the present.
\end{readersection}

\end{document}
